\section{Conclusion} \label{sec:Further_Steps}

This work aimed to develop a convolutional neural network that can classify exoplanet transit light curves while applying limited processing to the data. Other papers have applied heavy preprocessing as well as strong neural networks to achieve positive results. This paper showed an introduction to using simpler models with less provided information while maintaining results which are optimistic for future work.

First a binary model was developed which simply classified light curves between true exoplanet transits and eclipsing binaries. With the results from this model evaluation, the depth appeared to be a very important feature for the binary models classification and thus a more difficult 3-class CNN was developed. This model aimed to split the eclipsing binary dataset into shallow and deep groups based on the depth from the the normalized flux. 

The binary model returned a 92\% accuracy score indicating the majority of test data was correctly classified while the 88\% precision score warns that the model did classify a non-negligible batch of eclipsing binaries as transits. For a usable model this precision score is the most important to look into further. 

The multiclass model did not perform as well as the binary but the results did show the depth of transits was the major feature affecting the models decisions. As the results are macro-averaged, it is understood the shallow eclipsing binary class was the most challenging for the model. The CNN would mix up both transits and deep eclipsing binaries as the shallow case and this lowered all final metrics to the high 70s.

The main takeaway from this paper is the impact of depth for the classification problem. The results of both CNN models, between their confusion matrices and analysis of the misclassified instances, show the depth is correlated to the ability of the models to classify light curves correctly. 

The introduction of random forest models proved the CNN models are building off of the baseline. Both classifiers use different processes, with the random forest using inputted values while the CNN observes the light curve as a whole. It is clear the simple values such as depth, and flux ranges are important for classification. Potentially a future pipeline could combine calculated values and the CNN to make a more efficient combined model. 

This paper faced a few limitations which had to be followed. The first is the number of available confirmed exoplanet transits and valid eclipsing binaries. These were the only labelled Kepler objects that could be reliably used and in the future having more data to train and test off of will help later models. Without increasing the architecture of the CNN model, the classification was limited to 3 classes while there are several other potential false positive situations. 

The approaches in this paper work reasonably well for managing lightly processed Kepler data. As Kepler is a well documented mission, this makes for a great starting point for developing useful classification models. Shallow depth eclipsing binaries are the main obstacle for future models to handle.

\newpage
